\section{Determinant layers and the mesoscopic boundary}
\label{sec:determinant-boundary}

The large-sieve theorem becomes diagonal only after an $O(X)$ resolution
loss.  This section records exactly what survives below that scale.

For a sharp band put
\begin{equation}
 K_T(x):=
 \begin{cases}
  2\sin(Tx)/x,&x\ne0,\\
  2T,&x=0.
 \end{cases}
 \label{eq:sharp-band-kernel}
\end{equation}
Direct expansion of the original factor pairs gives
\begin{equation}
 \begin{aligned}
 \int_{-T}^{T}|\cR_{h,D}(X;t)|^2dt
 ={}&\sum_{\substack{n,m,n',m'\\n,n'>D}}
 u(n)u(n')\eta(mn+h)\eta(m'n'+h)\\
 &\quad\times W(mn/X)W(m'n'/X)
 K_T\!\left(\log\frac{nm'}{mn'}\right),
 \end{aligned}
 \label{eq:fourfold-hard-band}
\end{equation}
where $\eta(q)=\Lambda(q)-1$.  The highest-degree von Mangoldt term in
\eqref{eq:fourfold-hard-band} contains
\[
 \Lambda(n)\Lambda(n')\Lambda(mn+h)\Lambda(m'n'+h).
\]
Thus a small fixed-shift band is a genuine fourfold arithmetic
correlation, not merely a divisor mean value.

After ray compression, the same exact identity is
\begin{equation}
 \int_{-T}^{T}|\cR_{h,D}(X;t)|^2dt
 =\sum_{\substack{(a,b)=1\\(c,d)=1}}
 c_{a,b}\overline{c_{c,d}}
 K_T\!\left(\log\frac{ad}{bc}\right).
 \label{eq:ray-hard-kernel}
\end{equation}
The relevant near-resonance is measured by the integer
\begin{equation}
 \Delta((a,b),(c,d)):=ad-bc.
 \label{eq:ray-determinant}
\end{equation}

\begin{proposition}[Compact-Fourier determinant layers]
\label{prop:determinant-layers}
Let $\Phi\in L^1(\R)$ and suppose, with the Fourier convention used in
\eqref{eq:Fejer-transform}, that
$\supp\widehat\Phi\subset[-L,L]$.  Then
\begin{equation}
 \begin{aligned}
 \frac1T\int_\R\Phi(t/T)|\cR_{h,D}(X;t)|^2dt
 ={}&\sum_{\substack{(a,b)=1\\(c,d)=1}}
 c_{a,b}\overline{c_{c,d}}\\
 &\quad\times
 \widehat\Phi\!\left(T\log\frac{bc}{ad}\right).
 \end{aligned}
 \label{eq:compact-Fourier-layer-identity}
\end{equation}
If $L/T\leq1$, every nonzero summand satisfies
\begin{equation}
 |ad-bc|
 \leq2\beta X\sinh\!\left(\frac{L}{2T}\right)
 \ll_L\frac XT.
 \label{eq:determinant-layer-range}
\end{equation}
\end{proposition}

\begin{proof}
Insert \eqref{eq:ray-compression} and use Fourier inversion for the scaled
weight.  A term can survive only if
\[
 \left|\log\frac{ad}{bc}\right|\leq\frac LT.
\]
The exact hyperbolic identity from \cref{lem:ratio-spacing} gives
\[
 |ad-bc|
 =2\sqrt{abcd}\,
 \sinh\!\left(\frac12\left|\log\frac{ad}{bc}\right|\right).
\]
Since $ab,cd\leq\beta X$, this is at most the right-hand side of
\eqref{eq:determinant-layer-range}.
\end{proof}

At frequency scale $T=X^\theta$, the compact-Fourier mean square therefore
contains only determinant layers
\[
 |\Delta|\ll X^{1-\theta}.
\]
For $T\gg X$, only $\Delta=0$ remains; primitive fractions then force
$(a,b)=(c,d)$.  At $T\asymp X$, only finitely many nonzero determinant
layers can remain.  For $T=o(X)$, the permitted determinant range grows
and the number of potentially surviving layers is allowed to grow.  A
diagonal law then requires new cancellation among their arithmetic
coefficients.

\begin{proposition}[Sharpness of the $X$ threshold]
\label{prop:X-threshold-sharp}
Let $\Phi\in L^1(\R)$ have continuous Fourier transform with
$\widehat\Phi(0)\ne0$.  There are hyperbolic coefficient systems supported
on two primitive rays for which the $\Phi$-weighted form has a nonzero
off-diagonal term whenever $1\leq T\leq cX$, with $c>0$ depending only on the
kernel and support band.
Consequently the order $T\asymp X$ in
\cref{thm:exact-Fejer-ray-Parseval} cannot be uniformly improved for
arbitrary coefficients.
\end{proposition}

\begin{proof}
Take the two rays $(N,1)$ and $(N-1,1)$ with $N\asymp X$, and assign both
coefficient one.  Their determinant is $1$ and their frequency gap is
$\log(N/(N-1))=N^{-1}+O(N^{-2})$.  Continuity and
$\widehat\Phi(0)\ne0$ give an interval about zero on which
$\widehat\Phi$ is nonzero.  Choose $c$ so that every scaled gap with
$1\leq T\leq cX$ lies in that interval.  The cross term in
\eqref{eq:compact-Fourier-layer-identity} then survives.
\end{proof}

\subsection{The fixed-shift arithmetic left inside one ray}

The diagonal ray energy is not an elementary coefficient norm.  Expanding
one term gives
\begin{equation}
 \begin{aligned}
 |c_{a,b}|^2
 ={}&\sum_{k,\ell}
 u(ak)u(a\ell)
 \eta(abk^2+h)\eta(ab\ell^2+h)\\
 &\quad\times W(abk^2/X)W(ab\ell^2/X),
 \end{aligned}
 \label{eq:intra-ray-arithmetic}
\end{equation}
with the source cutoff understood.  Terms $k\ne\ell$ contain two target
prime weights on the quadratic sequences
\[
 abk^2+h,\qquad ab\ell^2+h.
\]
For example, the ray $(a,b)=(1,1)$ has coefficient
\begin{equation}
 c_{1,1}(h;X,D)
 =\sum_{k>D}u(k)\eta(k^2+h)W(k^2/X).
 \label{eq:quadratic-prime-ray}
\end{equation}
No fixed-$h$ asymptotic for \eqref{eq:quadratic-prime-ray} or for the
general off-diagonal \eqref{eq:intra-ray-arithmetic} is used or obtained.

The role of the long shift average is now transparent.  It does not make
the same-ray correlations disappear.  Rather, their total absolute
contribution is upper-bounded, while the $k=\ell$ target square acquires
the larger factor $\log H$.  Specializing the averaged theorem back to one
fixed shift would remove precisely that advantage.

\begin{remark}[The next analytic target]
A genuine advance into the mesoscopic region $1\ll T\ll X$ would estimate
the finitely many determinant-layer forms in
\eqref{eq:compact-Fourier-layer-identity} using their arithmetic
coefficients, rather than their absolute values.  Possible honest
weakenings include averaging $h$ over shorter windows, averaging the
determinant $\Delta$, or restricting $(a,b)$ to balanced subregions.  The
present paper supplies the exact layer identity, the permitted determinant
range, and the normalization against which such an estimate should be
measured.
\end{remark}
